\documentclass[11pt]{article}
\usepackage[margin=1in]{geometry}
\usepackage{amsmath, amssymb}
\usepackage[T1]{fontenc}
\usepackage{enumitem}
\usepackage{hyperref}
\usepackage{pifont}

\newcommand{\cmark}{\ding{51}}

\title{CS273P Final Exam 2026 Spring}
\date{}
\author{}

\begin{document}
\maketitle

\section*{Instructions}
Choose the best answer for each question.

\begin{enumerate}[leftmargin=*]

\item Suppose you run 1-nearest neighbor on a training set with no duplicate feature vectors and evaluate on the same training set. Which is true?

\begin{enumerate}[label=\Alph*.]
    \item Training error must be positive unless data are linearly separable
    \item Training error is always zero \cmark
    \item Training error depends on the number of classes
    \item Training error is zero only if features are normalized
\end{enumerate}


\item As $K$ increases in KNN classification, what usually happens?

\begin{enumerate}[label=\Alph*.]
    \item The model is more likely to overfit the training data
    \item The model is less likely to overfit, but may become more likely to underfit \cmark
    \item The model always fits the training data perfectly
    \item The model always becomes more accurate on test data
\end{enumerate}



\item Why is choosing $K$ using training accuracy alone usually a bad idea?

\begin{enumerate}[label=\Alph*.]
    \item Training accuracy cannot be computed for KNN
    \item Training accuracy tends to favor small $K$, often $K=1$ \cmark
    \item Training accuracy tends to favor large $K$
    \item Cross-validation is only for regression
\end{enumerate}

\item In high dimensions, which phenomenon most directly harms nearest-neighbor methods?

\begin{enumerate}[label=\Alph*.]
    \item The dot product becomes undefined
    \item Distances to nearest and farthest points become relatively similar \cmark
    \item Labels become continuous
    \item The number of classes must increase
\end{enumerate}

\item If every feature is multiplied by 1000, which classifier is most directly affected if no other preprocessing is done?

\begin{enumerate}[label=\Alph*.]
    \item Naive Bayes with Bernoulli features
    \item KNN with Euclidean distance \cmark
    \item PCA after centering and proper scaling
    \item Decision trees
\end{enumerate}

\item Which expression is correct?

\begin{enumerate}[label=\Alph*.]
    \item $P(y|x)=\frac{P(y)P(x)}{P(x|y)}$
    \item $P(y|x)=\frac{P(x|y)P(y)}{P(x)}$ \cmark
    \item $P(y|x)=P(x|y)P(x)$
    \item $P(y|x)=P(x|y)+P(y)$
\end{enumerate}

\item For classification, MAP prediction differs from maximum likelihood prediction because MAP includes:

\begin{enumerate}[label=\Alph*.]
    \item The evidence $P(x)$ only
    \item The prior $P(y)$ \cmark
    \item The covariance matrix only
    \item The learning rate
\end{enumerate}

\item Which statement is correct?

\begin{enumerate}[label=\Alph*.]
    \item $P(\theta|D)$ is the likelihood
    \item $P(D|\theta)$ is the posterior
    \item $P(D|\theta)$ is the likelihood \cmark
    \item $P(D)$ is the prior
\end{enumerate}

\item For binary features $x_j \in \{0,1\}$, the Bernoulli Naive Bayes class-conditional model typically writes
\[
P(x|y=c)=\prod_j p_{jc}^{x_j}(1-p_{jc})^{1-x_j}.
\]
What does $p_{jc}$ represent?

\begin{enumerate}[label=\Alph*.]
    \item The probability that class $c$ occurs
    \item The probability feature $j$ equals 1 given class $c$ \cmark
    \item The mean of feature $j$ over all data
    \item The posterior probability $P(y=c|x_j)$
\end{enumerate}

\item Which assumption is made by Gaussian Naive Bayes?

\begin{enumerate}[label=\Alph*.]
    \item Each class shares the same full covariance matrix
    \item Features are conditionally independent given the class, each modeled as Gaussian \cmark
    \item The posterior is Gaussian in the class label
    \item Data must be binary
\end{enumerate}

\item Let $x_1,\dots,x_n$ be i.i.d.\ from $\mathcal N(\mu,\sigma^2)$ with known $\sigma^2$. The MLE of $\mu$ is:

\begin{enumerate}[label=\Alph*.]
    \item $\sum_i x_i^2$
    \item $\frac{1}{n}\sum_i x_i$ \cmark
    \item $\frac{1}{n-1}\sum_i x_i$
    \item $\sqrt{\frac{1}{n}\sum_i x_i^2}$
\end{enumerate}

\item The Bayes optimal classifier assigns $x$ to the class:

\begin{enumerate}[label=\Alph*.]
    \item With largest prior probability
    \item With largest likelihood $P(x|y)$
    \item With largest posterior $P(y|x)$ \cmark
    \item With smallest entropy
\end{enumerate}

\item The standard empirical risk for linear regression with MSE over $n$ examples is:

\begin{enumerate}[label=\Alph*.]
    \item $\sum_{i=1}^n |w^T x_i-y_i|$
    \item $\frac{1}{n}\sum_{i=1}^n (w^T x_i-y_i)^2$ \cmark
    \item $\sum_{i=1}^n \log(1+e^{-y_i w^T x_i})$
    \item $\sum_{i=1}^n \max(0,1-y_iw^Tx_i)$
\end{enumerate}

\item Let
\[
L(w)=\frac{1}{n}\|Xw-y\|^2.
\]
What is $\nabla_w L(w)$?

\begin{enumerate}[label=\Alph*.]
    \item $\frac{1}{n}X^T(Xw-y)$
    \item $\frac{2}{n}X^T(Xw-y)$ \cmark
    \item $2X(Xw-y)$
    \item $Xw-y$
\end{enumerate}

\item Suppose $L(w)=\frac{1}{2}(wx-y)^2$, where $w,x,y$ are scalars. Starting at $w_0$, one gradient descent step with learning rate $\eta$ gives:

\begin{enumerate}[label=\Alph*.]
    \item $w_1 = w_0 - \eta(wx-y)$
    \item $w_1 = w_0 - \eta(w_0x-y)x$ \cmark
    \item $w_1 = w_0 + \eta(w_0x-y)x$
    \item $w_1 = w_0 - \eta(w_0-y)$
\end{enumerate}

\item Which statement is correct?

\begin{enumerate}[label=\Alph*.]
    \item SGD uses one example per update; batch GD uses all examples; mini-batch uses a subset \cmark
    \item SGD and mini-batch GD are identical
    \item Batch GD is always faster per epoch than mini-batch GD on GPUs
    \item Online GD means using all current examples seen so far in each step
\end{enumerate}

\item If gradient descent diverges during training of a smooth convex loss, the most likely explanation is:

\begin{enumerate}[label=\Alph*.]
    \item The learning rate is too small
    \item The learning rate is too large \cmark
    \item The dataset has too few examples
    \item The model is linear
\end{enumerate}

\item In binary logistic regression, $\sigma(w^T x)$ is interpreted as:

\begin{enumerate}[label=\Alph*.]
    \item The log-odds of class 1
    \item The probability of class 1 under the model \cmark
    \item The margin
    \item The hinge loss
\end{enumerate}

\item If $\sigma(z)=\frac{1}{1+e^{-z}}$, then $\sigma'(z)$ equals:

\begin{enumerate}[label=\Alph*.]
    \item $\sigma(z)$
    \item $1-\sigma(z)$
    \item $\sigma(z)(1-\sigma(z))$ \cmark
    \item $e^{-z}$
\end{enumerate}

\item For one example with label $y\in\{0,1\}$ and prediction $\hat y\in(0,1)$, the logistic loss is:

\begin{enumerate}[label=\Alph*.]
    \item $(y-\hat y)^2$
    \item $-y\log \hat y$
    \item $-[y\log \hat y+(1-y)\log(1-\hat y)]$ \cmark
    \item $\max(0,1-y\hat y)$
\end{enumerate}

\item For one example in logistic regression, with $\hat y=\sigma(w^T x)$, the gradient of the negative log-likelihood is:

\begin{enumerate}[label=\Alph*.]
    \item $(y-\hat y)x$
    \item $(\hat y-y)x$ \cmark
    \item $(1-\hat y)x$
    \item $-yx$
\end{enumerate}

\item The decision boundary of binary logistic regression with threshold 0.5 is:

\begin{enumerate}[label=\Alph*.]
    \item Quadratic in the features
    \item Linear in the features \cmark
    \item Always nonlinear because sigmoid is nonlinear
    \item Undefined unless features are normalized
\end{enumerate}

\item For logits $z\in\mathbb R^C$, softmax probabilities satisfy:

\begin{enumerate}[label=\Alph*.]
    \item Each output is between $-1$ and 1
    \item Outputs sum to 1 \cmark
    \item Outputs are binary
    \item Outputs are independent of additive shifts to logits only if $C=2$
\end{enumerate}

\item Which transformation leaves softmax probabilities unchanged?

\begin{enumerate}[label=\Alph*.]
    \item Multiplying all logits by 2
    \item Adding the same constant to all logits \cmark
    \item Squaring all logits
    \item Replacing logits with their signs
\end{enumerate}

\item If $y$ is one-hot and $\hat y$ is the softmax output, then
\[
-\sum_i y_i \log \hat y_i
\]
equals:

\begin{enumerate}[label=\Alph*.]
    \item The negative log of the predicted probability of the true class \cmark
    \item The sum of all logits
    \item The mean squared error
    \item The entropy of the input features
\end{enumerate}

\item For softmax with cross-entropy loss for one example, the gradient with respect to the logits is:

\begin{enumerate}[label=\Alph*.]
    \item $\hat y-y$ \cmark
    \item $y-\hat y$
    \item $(\hat y-y)^2$
    \item $\hat y(1-\hat y)$
\end{enumerate}

\item Adding $\lambda \|w\|_2^2$ to the loss tends to:

\begin{enumerate}[label=\Alph*.]
    \item Encourage many weights to be exactly zero more strongly than L1
    \item Penalize large weights smoothly \cmark
    \item Make the optimization non-differentiable everywhere
    \item Leave the optimum unchanged
\end{enumerate}

\item A key reason L1 regularization may be preferred over L2 is that L1 often:

\begin{enumerate}[label=\Alph*.]
    \item Gives closed-form solutions for deep networks
    \item Encourages sparse solutions \cmark
    \item Avoids overfitting in every setting
    \item Is equivalent to dropout
\end{enumerate}

\item For objective
\[
J(w)=\frac{1}{n}\|Xw-y\|^2 + \lambda \|w\|_2^2,
\]
the gradient is:

\begin{enumerate}[label=\Alph*.]
    \item $\frac{2}{n}X^T(Xw-y)+2\lambda w$ \cmark
    \item $\frac{2}{n}X(Xw-y)+\lambda$
    \item $X^T(Xw-y)+\lambda |w|$
    \item $\frac{1}{n}(Xw-y)^2+\lambda w$
\end{enumerate}

\item In a linear hard-margin SVM, the geometric margin is increased by:

\begin{enumerate}[label=\Alph*.]
    \item Increasing $\|w\|$
    \item Decreasing $\|w\|$ while satisfying separation constraints \cmark
    \item Making all slack variables large
    \item Using sigmoid instead of hinge loss
\end{enumerate}

\item In the soft-margin SVM objective
\[
\frac{1}{2}\|w\|^2 + C\sum_i \xi_i,
\]
a larger $C$ generally means:

\begin{enumerate}[label=\Alph*.]
    \item Stronger penalty for violations of the margin constraints \cmark
    \item Weaker penalty for violations
    \item No regularization
    \item Larger learning rate
\end{enumerate}

\item For hinge loss $L=\max(0,1-yf(x))$, the loss is zero when:

\begin{enumerate}[label=\Alph*.]
    \item $yf(x)\le 0$
    \item $yf(x)\ge 1$ \cmark
    \item $f(x)=0$
    \item $y=0$
\end{enumerate}

\item For $L(w)=\max(0,1-yw^Tx)$, which is a valid subgradient (derivative) of the loss with respect to $w$ when $yw^Tx < 1$?

\begin{enumerate}[label=\Alph*.]
    \item $0$
    \item $yx$
    \item $-yx$ \cmark
    \item $-x$
\end{enumerate}

\item Which statement is correct?

\begin{enumerate}[label=\Alph*.]
    \item Hinge loss is smooth everywhere, logistic loss is not
    \item Logistic loss gives zero penalty once margin exceeds 1
    \item Hinge loss is piecewise linear and not differentiable at the margin \cmark
    \item Both losses are identical up to scaling
\end{enumerate}

\item Why does adding nonlinear activation functions between layers matter?

\begin{enumerate}[label=\Alph*.]
    \item Without them, a multi-layer network collapses to an equivalent linear map \cmark
    \item Without them, gradient descent cannot run
    \item Without them, backpropagation is impossible
    \item Without them, the model automatically becomes recurrent
\end{enumerate}

\item Backpropagation is best described as:

\begin{enumerate}[label=\Alph*.]
    \item Forward computation of activations only
    \item Efficient application of the chain rule from output layer to earlier layers \cmark
    \item Random search over parameters
    \item A method that only works for sigmoid networks
\end{enumerate}

\item Suppose \texttt{x} has shape \texttt{(batch\_size, 20)}, first layer is \texttt{nn.Linear(20, 50)}, second layer is \texttt{nn.Linear(50, 3)}. What is the output shape?

\begin{enumerate}[label=\Alph*.]
    \item \texttt{(20, 3)}
    \item \texttt{(batch\_size, 50)}
    \item \texttt{(batch\_size, 3)} \cmark
    \item \texttt{(3, batch\_size)}
\end{enumerate}

\item Which is true about ReLU $f(z)=\max(0,z)$?

\begin{enumerate}[label=\Alph*.]
    \item It saturates for large positive $z$
    \item It often helps mitigate vanishing gradients compared with sigmoid \cmark
    \item Its output is always in $(0,1)$
    \item It is differentiable everywhere
\end{enumerate}

\item For a 2D convolution with input size $32\times 32$, kernel size $5$, stride $2$, padding $1$, the output spatial size is:

\begin{enumerate}[label=\Alph*.]
    \item $15\times 15$ \cmark
    \item $16\times 16$
    \item $17\times 17$
    \item $14\times 14$
\end{enumerate}

\item Which line defines a convolution mapping of images with 3 input channels to 8 output channels with kernel size 3 and padding 1?

\begin{enumerate}[label=\Alph*.]
    \item \texttt{nn.Conv2d(3, 8, kernel\_size=3, padding=1)} \cmark
    \item \texttt{nn.Conv2d(8, 3, kernel\_size=3, padding=1)}
    \item \texttt{nn.Linear(3, 8, 3)}
    \item \texttt{nn.Conv1d(3, 8, kernel\_size=3, padding=1)}
\end{enumerate}

\item In a standard convolution layer, the number of output channels is determined by:

\begin{enumerate}[label=\Alph*.]
    \item The stride
    \item The padding
    \item The number of filters \cmark
    \item The input width
\end{enumerate}

\item Why do vanilla RNNs struggle with long-range dependencies?

\begin{enumerate}[label=\Alph*.]
    \item They cannot process variable-length sequences
    \item Repeated multiplication through time can cause gradients to shrink or explode \cmark
    \item They do not use matrix multiplication
    \item They require one-hot inputs only
\end{enumerate}

\item LSTM cells improve vanilla RNNs primarily by:

\begin{enumerate}[label=\Alph*.]
    \item Replacing matrix multiplication with convolution
    \item Using gating mechanisms to better preserve or forget information over time \cmark
    \item Eliminating recurrence entirely
    \item Forcing hidden states to be binary
\end{enumerate}

\item A sequence-to-sequence model is most naturally suited for:

\begin{enumerate}[label=\Alph*.]
    \item Binary image classification only
    \item Mapping an input sequence to an output sequence, possibly of different length \cmark
    \item Clustering fixed-dimensional points only
    \item PCA on tabular data
\end{enumerate}

\item For sentiment analysis on sentences, which model type is especially natural because it handles variable-length token sequences?

\begin{enumerate}[label=\Alph*.]
    \item K-means
    \item RNN/LSTM or Transformer \cmark
    \item Linear regression
    \item PCA
\end{enumerate}

\item In scaled dot-product attention, the attention weights are computed as:

\begin{enumerate}[label=\Alph*.]
    \item $\text{softmax}(Q+K)V$
    \item $\text{softmax}(QK^T/\sqrt{d_k})V$ \cmark
    \item $Q(K^TV)$
    \item $\text{softmax}(QV^T/\sqrt{d_k})K$
\end{enumerate}

\item In transformer attention, dividing by $\sqrt{d_k}$ helps because otherwise for large $d_k$:

\begin{enumerate}[label=\Alph*.]
    \item Dot products can become large, making softmax too peaky and gradients unstable \cmark
    \item The values matrix $V$ becomes singular
    \item The sequence length changes
    \item The model becomes recurrent
\end{enumerate}

\item K-means clustering aims to minimize:

\begin{enumerate}[label=\Alph*.]
    \item The sum of distances between all pairs of points
    \item The within-cluster sum of squared distances to cluster centroids \cmark
    \item The determinant of the covariance matrix
    \item The classification error rate
\end{enumerate}

\item Which statement best distinguishes Gaussian mixture models from K-means?

\begin{enumerate}[label=\Alph*.]
    \item GMM gives only hard assignments; K-means gives soft assignments
    \item GMM models clusters probabilistically and can provide soft assignments \cmark
    \item K-means models full covariance Gaussians
    \item They are identical if initialized the same way
\end{enumerate}

\item Which statement is correct?

\begin{enumerate}[label=\Alph*.]
    \item PCA always gives nonlinear dimension reduction
    \item A denoising autoencoder is trained to reconstruct the clean input from a corrupted version \cmark
    \item A VAE is trained only with reconstruction loss and no regularization term
    \item Standard autoencoders require class labels
\end{enumerate}


\item In softmax multiclass classification, let $z \in \mathbb{R}^C$ be the logits, $\hat y_i = \frac{e^{z_i}}{\sum_j e^{z_j}}$, and let the target $y$ be one-hot. Which statement is correct?

\begin{enumerate}[label=\Alph*.]
    \item Adding a constant to every coordinate of $z$ changes the loss but not the predicted class
    \item Adding a constant to every coordinate of $z$ leaves both the softmax probabilities and cross-entropy loss unchanged \cmark
    \item Multiplying every coordinate of $z$ by 2 leaves the softmax probabilities unchanged
    \item The cross-entropy loss depends only on the largest logit
\end{enumerate}

\item Consider PCA on centered data matrix $X \in \mathbb{R}^{n\times d}$. Suppose $v_1$ is the first principal direction, constrained by $\|v_1\|=1$. Which optimization problem does $v_1$ solve?

\begin{enumerate}[label=\Alph*.]
    \item \(\displaystyle \min_{\|v\|=1} \|X - Xvv^T\|_F^2\) and equivalently \(\max_{\|v\|=1} v^T X^T X v\) \cmark
    \item \(\displaystyle \max_{\|v\|=1} \|Xv\|_1\)
    \item \(\displaystyle \min_{\|v\|=1} v^T X^T X v\)
    \item \(\displaystyle \max_{\|v\|=1} \det(X^T X)\)
\end{enumerate}

\item In a linear SVM with hinge loss and $L_2$ regularization,
\[
J(w)=\frac{\lambda}{2}\|w\|^2 + \sum_{i=1}^n \max(0,1-y_i w^T x_i),
\]
where $y_i \in \{-1,+1\}$. Which is a valid subgradient of $J(w)$?

\begin{enumerate}[label=\Alph*.]
    \item \(\displaystyle \lambda w - \sum_{i:\, y_i w^T x_i < 1} y_i x_i\) \cmark
    \item \(\displaystyle \lambda w + \sum_{i=1}^n y_i x_i\)
    \item \(\displaystyle \sum_{i:\, y_i w^T x_i > 1} y_i x_i\)
    \item \(\displaystyle -\lambda w - \sum_{i:\, y_i w^T x_i < 1} y_i x_i\)
\end{enumerate}

\item In scaled dot-product attention, let $Q \in \mathbb{R}^{m \times d_k}$, $K \in \mathbb{R}^{n \times d_k}$, and $V \in \mathbb{R}^{n \times d_v}$. The attention output is
\[
\mathrm{Attention}(Q,K,V)=\mathrm{softmax}\left(\frac{QK^T}{\sqrt{d_k}}\right)V.
\]
What is the shape of the output?

\begin{enumerate}[label=\Alph*.]
    \item $m \times n$
    \item $n \times d_v$
    \item $m \times d_v$ \cmark
    \item $d_k \times d_v$
\end{enumerate}


\item Suppose a 2D convolution layer has input tensor shape \texttt{(batch\_size, 16, 31, 31)} in PyTorch format \texttt{(N,C,H,W)}. The layer is
\texttt{nn.Conv2d(16, 32, kernel\_size=5, stride=2, padding=2)}.
What is the output shape?

\begin{enumerate}[label=\Alph*.]
    \item \texttt{(batch\_size, 32, 16, 16)} \cmark
    \item \texttt{(batch\_size, 32, 15, 15)}
    \item \texttt{(batch\_size, 16, 16, 16)}
    \item \texttt{(batch\_size, 32, 31, 31)}
\end{enumerate}

\item Consider a vanilla RNN
\[
h_t = \tanh(W_{hh} h_{t-1} + W_{xh} x_t + b).
\]
Which statement best explains why vanishing gradients can occur when backpropagating through many time steps?

\begin{enumerate}[label=\Alph*.]
    \item Because $\tanh$ is not differentiable at 0
    \item Because repeated multiplication by Jacobians with singular values less than 1 can make gradients decay exponentially \cmark
    \item Because the hidden state has fixed dimension
    \item Because the input sequence is too short
\end{enumerate}

\item In an LSTM, which component most directly allows the cell state to preserve information over long time spans with relatively stable gradient flow?

\begin{enumerate}[label=\Alph*.]
    \item The output gate alone
    \item The additive cell-state update controlled by forget and input gates \cmark
    \item The softmax layer on top of the network
    \item The hidden state dimension being larger than the input dimension
\end{enumerate}

\item In multi-head self-attention, why is it useful to project the input into multiple different query/key/value subspaces instead of using a single attention head?

\begin{enumerate}[label=\Alph*.]
    \item It guarantees fewer parameters than single-head attention
    \item It allows the model to attend to different types of relationships or positions simultaneously \cmark
    \item It removes the need for positional encodings
    \item It makes the attention computation linear in sequence length
\end{enumerate}

\item Suppose self-attention is applied to a sequence of length $n$ with hidden dimension $d$. Ignoring constant factors and assuming $d$ is fixed across layers, the dominant time complexity of forming the attention matrix in standard self-attention is:

\begin{enumerate}[label=\Alph*.]
    \item $O(n)$
    \item $O(n \log n)$
    \item $O(n d)$
    \item $O(n^2 d)$ \cmark
\end{enumerate}

\item Let
\[
L(w)=\frac{1}{2}(w^T x-y)^2
\]
for one training example, where $w,x\in\mathbb{R}^d$. What is the Hessian $\nabla^2_w L(w)$?

\begin{enumerate}[label=\Alph*.]
    \item $xx^T$ \cmark
    \item $(w^Tx-y)xx^T$
    \item $2xx^T$
    \item $(w^Tx-y)x$
\end{enumerate}

\item For binary logistic regression with data matrix $X\in\mathbb{R}^{n\times d}$, predictions $p_i=\sigma(w^T x_i)$, and negative log-likelihood
\[
L(w)=-\sum_{i=1}^n \left[y_i\log p_i+(1-y_i)\log(1-p_i)\right],
\]
which statement is true about $L(w)$?

\begin{enumerate}[label=\Alph*.]
    \item $L(w)$ is generally non-convex because of the sigmoid
    \item $L(w)$ is convex because its Hessian can be written as $X^T D X$ with $D$ diagonal and nonnegative \cmark
    \item $L(w)$ is concave when the data are linearly separable
    \item $L(w)$ is convex only when $d=1$
\end{enumerate}

\item Let $X$ be a centered data matrix. In PCA, the total variance explained by the first $k$ principal components is:

\begin{enumerate}[label=\Alph*.]
    \item The sum of the first $k$ singular values divided by the largest singular value
    \item The sum of the first $k$ eigenvalues of the covariance matrix divided by the sum of all eigenvalues \cmark
    \item The product of the first $k$ eigenvalues divided by the trace
    \item The trace of $X^T X$ divided by $k$
\end{enumerate}

\item In a variational autoencoder, the loss for one example is typically written as
\[
\mathcal{L}(x)= -\mathbb{E}_{q_\phi(z|x)}[\log p_\theta(x|z)] + \mathrm{KL}(q_\phi(z|x)\,\|\,p(z)).
\]
What is the main role of the KL term?

\begin{enumerate}[label=\Alph*.]
    \item To force the decoder to be linear
    \item To encourage the approximate posterior to stay close to the prior \cmark
    \item To maximize reconstruction error
    \item To make latent variables deterministic
\end{enumerate}

\item In $k$-means clustering, suppose cluster assignments are fixed. What is the optimal update for each cluster center?

\begin{enumerate}[label=\Alph*.]
    \item The point in the cluster closest to the origin
    \item The coordinate-wise median of the assigned points
    \item The mean of the assigned points \cmark
    \item The point in the cluster with largest norm
\end{enumerate}

\item In a Gaussian mixture model with parameters $\theta = \{\pi_k, \mu_k, \Sigma_k\}_{k=1}^K$, the marginal log-likelihood for one data point $x$ is
\[
\log p(x) = \log \sum_{k=1}^K \pi_k \, \mathcal{N}(x \mid \mu_k, \Sigma_k).
\]
Why is the EM algorithm useful for maximizing this objective?

\begin{enumerate}[label=\Alph*.]
    \item Because the logarithm turns the sum over mixture components into a product
    \item Because introducing latent component indicators makes the complete-data log-likelihood easier to optimize \cmark
    \item Because EM always finds the global maximum of the likelihood
    \item Because the Gaussian covariance matrices do not need to be estimated
\end{enumerate}

\item Let $q(z)$ be any distribution over the latent variable $z$ in a latent-variable model. The evidence lower bound (ELBO) for $\log p(x)$ is
\[
\mathrm{ELBO}(q) = \mathbb{E}_{q(z)}[\log p(x,z)] - \mathbb{E}_{q(z)}[\log q(z)].
\]
Which identity is correct?

\begin{enumerate}[label=\Alph*.]
    \item $\log p(x) = \mathrm{ELBO}(q) - \mathrm{KL}(q(z)\|p(z|x))$
    \item $\log p(x) = \mathrm{ELBO}(q) + \mathrm{KL}(q(z)\|p(z|x))$ \cmark
    \item $\log p(x) = \mathrm{ELBO}(q) + \mathrm{KL}(p(z|x)\|q(z))$
    \item $\log p(x) = \mathrm{KL}(q(z)\|p(z|x)) - \mathrm{ELBO}(q)$
\end{enumerate}

\item Consider a two-layer neural network for one training example:
\[
z_1 = W_1 x + b_1,\qquad a_1 = \mathrm{ReLU}(z_1),
\]
\[
z_2 = W_2 a_1 + b_2,\qquad \hat y = z_2,
\]
and the loss is
\[
L = \frac{1}{2}\|\hat y - y\|_2^2.
\]
Which expression is correct for the gradient of the loss with respect to $W_1$? ($\odot$ is elementwise product, $\mathbf{1}_{z_1>0}$ is an indicator function that is 1 if $z_1>0$ and 0 otherwise)

\begin{enumerate}[label=\Alph*.]
    \item \(\displaystyle \frac{\partial L}{\partial W_1} = (\hat y - y)x^T\)
    \item \(\displaystyle \frac{\partial L}{\partial W_1} = \bigl(W_2^T(\hat y-y)\odot \mathbf{1}_{z_1>0}\bigr)x^T\) \cmark
    \item \(\displaystyle \frac{\partial L}{\partial W_1} = W_1^T(\hat y-y)\)
    \item \(\displaystyle \frac{\partial L}{\partial W_1} = \bigl((\hat y-y)\odot \mathbf{1}_{z_1>0}\bigr)x^T\)
\end{enumerate}


\end{enumerate}
\end{document}
